\chapter{Verification Approach and Implementation}\label{verification-approach-and-implementation}

\section{Architecture}\label{architecture}

Figure \ref{fig:evidence-first-architecture} shows the implemented system as a
UML-style component architecture. Dashed subsystem boundaries separate input
preparation, evidence reasoning, and decision and trace generation. The
solid connectors denote verdict-bearing component dependencies. Dashed
connectors distinguish the localization trace: the region grounder is a
first-class component, but its coordinates are carried through aggregation as
evidence metadata rather than used to determine the requirement label.

\begin{figure}[htbp]
  \centering
  \resizebox{\textwidth}{!}{\begin{tikzpicture}[
  component/.style={
    draw=TUMDarkGray,
    rounded corners=1.5pt,
    minimum height=0.95cm,
    align=center,
    inner xsep=7pt,
    inner ysep=4pt,
    font=\small
  },
  external/.style={
    component,
    fill=TUMLightGray!18
  },
  subsystem/.style={
    draw=TUMGray,
    dashed,
    rounded corners=5pt,
    inner sep=8pt
  },
  dependency/.style={->, >=latex, draw=TUMDarkGray, line width=0.7pt},
  trace/.style={->, >=latex, draw=TUMGray, line width=0.7pt, dashed},
  connectorlabel/.style={font=\scriptsize, fill=\bg, inner sep=1.2pt, text=TUMDarkGray}
]

\node[external, text width=2.25cm] (screens) at (0,1.35) {
  \textbf{Screenshot flow}\\[-1pt]
  \scriptsize ordered PNG files
};
\node[external, text width=2.25cm] (requirements) at (0,-1.35) {
  \textbf{Requirements}\\[-1pt]
  \scriptsize reviewed or raw text
};

\node[component, text width=2.55cm] (ingestion) at (3.65,1.35) {
  \scriptsize\guillemotleft component\guillemotright\\[-1pt]
  \textbf{Flow Ingestion}
};
\node[component, text width=2.55cm] (understanding) at (3.65,-1.35) {
  \scriptsize\guillemotleft component\guillemotright\\[-1pt]
  \textbf{Requirement}\\[-2pt]
  \textbf{Understanding}
};
\node[subsystem, fit=(ingestion)(understanding),
  label={[font=\small\bfseries, anchor=south west]north west:Input preparation}] (preparation) {};

\node[component, text width=2.55cm] (retrieval) at (7.45,1.35) {
  \scriptsize\guillemotleft component\guillemotright\\[-1pt]
  \textbf{Evidence Retriever}
};
\node[component, text width=2.55cm] (verifier) at (7.45,-1.35) {
  \scriptsize\guillemotleft component\guillemotright\\[-1pt]
  \textbf{Multimodal}\\[-2pt]
  \textbf{Verifier}
};
\node[subsystem, fit=(retrieval)(verifier),
  label={[font=\small\bfseries, anchor=south west]north west:Evidence reasoning}] (reasoning) {};

\node[component, text width=2.55cm] (aggregation) at (11.25,1.35) {
  \scriptsize\guillemotleft component\guillemotright\\[-1pt]
  \textbf{Label Aggregator}
};
\node[component, text width=2.55cm] (grounding) at (11.25,-1.35) {
  \scriptsize\guillemotleft component\guillemotright\\[-1pt]
  \textbf{Region Grounder}
};
\node[subsystem, fit=(aggregation)(grounding),
  label={[font=\small\bfseries, anchor=south west]north west:Decision and trace}] (decision) {};

\node[external, text width=2.25cm] (result) at (14.85,1.35) {
  \textbf{Verification result}\\[-1pt]
  \scriptsize label, evidence, usage
};

\draw[dependency] (screens) -- node[connectorlabel, above] {ordered flow} (ingestion);
\draw[dependency] (requirements) -- node[connectorlabel, above] {requirement} (understanding);
\draw[dependency] (ingestion) -- (retrieval);
\draw[dependency] (understanding.north east) to[out=35,in=215] (retrieval.south west);
\draw[dependency] (retrieval.south) -- node[connectorlabel, right] {evidence} (verifier.north);
\draw[dependency] (understanding.east) -- (verifier.west);
\draw[dependency] (verifier.north east) -- (aggregation.south west);
\draw[trace] (verifier.south east) to[out=-25,in=205]
  node[connectorlabel, below] {regions} (grounding.south west);
\draw[trace] (grounding.north) -- (aggregation.south);
\draw[dependency] (aggregation) -- node[connectorlabel, above] {result} (result);

\end{tikzpicture}
}
  \caption{UML-style component view of the evidence-first verifier. Dashed
  boundaries group related components; solid connectors show verdict-bearing
  dependencies, while dashed connectors carry region-level trace metadata.}
  \label{fig:evidence-first-architecture}
\end{figure}

Each component consumes and emits typed structured records rather than an
unstructured chat transcript. Inputs preserve flow identifiers, step indices,
screenshots, and requirements. Outputs preserve labels, claims, evidence
units, uncertainty reasons, rationales, model metadata, and usage information.
This record contract is already represented by the component dependencies in
Figure \ref{fig:evidence-first-architecture}; a second linear data-flow diagram
would repeat the same transformation at a different level of abstraction.

Dependency injection allows the deterministic and multimodal components to be
compared without changing this data contract.

The staged architecture exposes four failure locations: claim decomposition,
screenshot retrieval, visual interpretation, and final aggregation. Its
purpose is diagnostic. Accuracy relative to a direct prompt is measured in the
experiments rather than assumed from the architecture.

\section{Flow Ingestion and Screen Representation}\label{flow-ingestion-and-screen-representation}

A flow directory contains a stable flow identifier, ordered screenshot files,
task metadata, and step metadata. The ingestion layer preserves the recorded
step index and image dimensions. Original-resolution screenshots are retained
locally alongside processed assets so that evaluation can distinguish model
input resolution from review resolution.

Screen understanding constructs a lightweight representation for each step. In
the main benchmark, it records image dimensions and extracts visible-text
candidates from the available Mind2Web HTML metadata. Existing OCR sidecars are
added where available. The resulting text supports inexpensive retrieval, while
the screenshot remains the authoritative visual input to the multimodal
verifier. A text match is never itself interpreted as proof of fulfillment.

Local flow material is separated from versioned annotations. A fresh checkout
contains code, requirements, manifests, and evaluation configuration, while the
Mind2Web source data must be obtained separately and exported into the expected
local structure. This repository layout implements the redistribution boundary
described in Section 7.7.

\section{Requirement Understanding}\label{requirement-understanding}

Requirement understanding performs two optional operations: UI-evaluability
classification and claim decomposition. The deterministic classifier uses
rules over the requirement text and serves as a low-cost baseline. The main
multimodal runs can instead request an evaluability judgment from the model.
Neither prediction is silently substituted for the reviewed reference during
evaluation.

The raw claim policy treats the complete requirement as one verification unit.
The gated policy applies deterministic decomposition only when the requirement
contains structures that plausibly benefit from splitting. In the final
benchmark it produces 281 claims from 258 requirements and splits 31
requirements. An LLM decomposition fallback exists in the general application
but is disabled in the controlled matrix so that hidden additional model calls
do not confound the comparison.

The repository also contains 541 reviewed benchmark claims. Supplying these
claims to the verifier creates an oracle or provided-claim condition. It is
useful for diagnosing claim-status prediction when the decomposition is fixed,
but it must not be reported as the performance of automatic decomposition.

\section{Screenshot Selection}\label{screenshot-selection}

The retriever scores screenshot steps for each claim using the textual screen
representations. The default controlled policy is lexical retrieval, chosen
because it is deterministic, inexpensive, and easy to reproduce. Alternative
implementations support TF-IDF, local embeddings, and text-only LLM reranking,
but they are not mixed into the primary matrix.

For the top-4 condition, lexical retrieval first produces item-level candidate
steps. For each verification batch of up to eight items, these candidates are
then compressed into a shared set of at most four attached screenshots. The
condition therefore evaluates retrieval together with batch-level evidence
compression; it does not guarantee four independently selected screenshots for
each requirement or claim.

The evaluated flows do not require retrieval to fit within the model input
limit. The longest flow contains 19 screenshots and remains well within this
limit. Shared top-k selection instead places a fixed bound on the visual input and
produces an explicit evidence ranking. The comparison tests whether this
smaller evidence set preserves labels and traces at lower cost without omitting
decisive screenshots. These properties may become more important for longer
flows or models with smaller context windows.

The capacity check used during system design reaches the same conclusion. Even
a conservative 1,000-line, 20-screenshot stress case fits within the documented
native context of Qwen3-VL-8B-Instruct and the input limits of Gemini 2.5
Flash-Lite and Gemini 3.1 Flash-Lite \autocite{qwen3vl8b2025,googlegemini25flashlite2026,googlegemini31flashlite2026}. Retrieval is
consequently evaluated as a bounded evidence and efficiency policy, not as a
necessary workaround for the present benchmark. Section 6.3 reports whether
that policy preserves decisive evidence and reduces end-to-end cost.

The whole-flow condition attaches every screenshot in its recorded order. It
avoids retrieval omission but increases image input and irrelevant context. The
late-state failures are especially relevant because lexical overlap may favor
an early input form over a later result or summary screen.

A separate order-unavailable robustness condition keeps the screenshot set,
requirements, model, chunking, label contract, and aggregation matched to the
raw whole-flow condition. For each flow, a fixed recorded permutation changes
the attachment order and replaces original step identifiers with local apparent
identifiers. The prompt discloses that the original chronology was removed by
the test environment and instructs the model not to invent before/after
relations. OCR and image metadata remain attached to the corresponding image,
and cited apparent identifiers are mapped back to original flow steps before
evidence scoring. This manipulation tests verification when trustworthy temporal
order is unavailable; it does not deceive the model with a false chronology.

\section{Multimodal Claim Verification}\label{multimodal-claim-verification}

The multimodal verifier receives the fixed label definitions, the requirement
or claims, ordered screenshots, and stable step identifiers. It returns
structured JSON containing UI evaluability where requested, claim statuses,
evidence steps, uncertainty reasons, rationales, and the fields needed for
requirement-level aggregation. Responses are validated before they enter the
evaluation. Missing or malformed items remain coverage failures rather than
being converted to semantic abstentions.

The final controlled matrix uses \texttt{gemini-3.1-flash-lite} with temperature zero,
thinking level \texttt{low}, fixed prompt and aggregation versions, and deterministic
chunks of at most eight requirements. Gemini 2.5 Flash-Lite provides a matched
low-cost model comparison with thinking budget zero. Qwen3-VL-8B-Instruct,
served through OpenRouter with provider fallbacks disabled, provides the hosted
open-weight baseline. Exact provider, model, parameters, execution date,
tokens, failures, and costs are archived following the reporting guidelines
for empirical software-engineering studies involving LLMs proposed by \textcite{baltes2026}.

Gemini 3.1 Flash-Lite was selected as the primary experimental model because it
supports multi-image input and structured output, was available through the
existing verified adapter, and was inexpensive enough to run the complete
matrix and repetitions while holding the model fixed. Gemini 2.5 Flash-Lite was
selected as a deliberately cheaper earlier-generation sensitivity baseline
rather than a claim about frontier performance. The hosted Qwen model addresses
the methodological value of an open-weight comparison without implying full
local reproducibility. A local SmolVLM2 pilot was documented but not promoted
to a benchmark result because inference on the available hardware was too slow
and the capped response was incomplete.

The historical Gemini 3.1 Pro preview run provides contextual strong-model
evidence. Its preview status, cost, and earlier prompt and grouping
configuration prevent its use in the controlled matrix. The model set was
chosen for matched comparisons rather than leaderboard coverage.

Temperature zero is not treated as a guarantee of determinism. The central
Gemini and Qwen configurations were therefore repeated. Each repetition uses a
separate output and cache directory, and the stability analysis distinguishes
label invariance from variation in cited evidence steps.

\section{Label Aggregation}\label{label-aggregation}

The label aggregator is deterministic and shared across the controlled
configurations. It receives claim statuses, claim importance, evidence
presence, uncertainty reasons, and effective UI evaluability. It enforces the
central safety gates: no \texttt{FULFILLED} label without evidence, no
\texttt{NOT\_FULFILLED} label without a visible contradiction, and no fulfilled
decision while a material hidden property remains unresolved.

Deterministic aggregation reduces prompt dependence at the requirement-label
boundary. The model still determines the semantic claim statuses, and poor
upstream decisions remain unrecoverable. The weak deterministic baseline
demonstrates this limitation.

The offline forced-decision analysis changes only the final policy applied to
frozen model outputs. It maps native abstentions to \texttt{NOT\_FULFILLED} without
performing new inference. Because it cannot create a positive label, it is
interpreted as a closed-world counterfactual rather than a test of whether the
model learned to abstain.

\section{Region-Level Evidence Grounding}\label{region-level-evidence-grounding}

Region grounding follows screenshot selection and semantic verification. Early
variants asked the multimodal model to emit normalized coordinates directly.
These boxes were unstable and sometimes landed on nearby text rather than the
specific evidence. OCR refinement improved text-aligned cases but could not
represent non-text controls or repair an incorrect semantic description.

The later candidate-mark pipeline generates candidate regions from OCR and UI
detectors, overlays stable identifiers, asks the multimodal model to select the
minimal evidence marks, and maps the selected identifiers back to pixel
coordinates deterministically. A grid can serve as a fallback when no suitable
candidate exists. Candidate generation and mark selection are separate sources
of error: proposal coverage limits the best achievable grounding result, while
the verifier may still choose the wrong available proposal.

The output permits zero, one, or several boxes per claim and supports an
explicit \texttt{NO\_\allowbreak VISIBLE\_\allowbreak REGION} response. Coordinates are stored together with the
image path, width, height, asset hash, coordinate space, and raw and refined
values where applicable. The review interface rescales boxes independently when
it displays a higher-resolution copy of the screenshot.

Bounding boxes provide a more precise evidence trace. Their evaluation measures
geometric validity, relevance, and sufficiency. Label accuracy remains outside
this comparison because region generation is not used as an accuracy
intervention.

\section{Annotation and Review Workbench}\label{annotation-and-review-workbench}

The project includes a local web workbench for browsing flows, reviewing
candidate and gold requirements, inspecting claims and evidence, running ad-hoc
verification, and comparing predicted regions with review judgments. Review
responses are stored separately from production outputs so that inspecting a
prediction does not silently overwrite the benchmark.

The UI-evaluability audit presents every case in which the stored author label
differs from the deterministic classifier. It shows both labels, the
classifier's triggered hidden-property terms, the author annotation note, claim
composition, and a generated diagnostic hypothesis. The author then records a
resolution and whether a later reference amendment is recommended. The region
audit has a different purpose: it shows the frozen V7 regions and asks whether
they are geometrically valid, semantically relevant, and sufficient for the
stated claim.

Both audits are performed by the thesis author. Their samples target observed
disagreements and predicted regions, so their results are reported as
qualitative boundary and output analyses. These diagnostic audits are separate
from the prediction-independent reference-annotation workflow and do not
define the labels used in the headline comparisons.

\section{Reproducibility and Artifact Management}\label{reproducibility-and-artifact-management}

Experiment configurations specify model identifiers, claim and screenshot
policies, retrieval parameters, generation settings, cost guards, and expected
coverage. Preflight manifests record the Git state, environment versions,
benchmark counts, gold hashes, prompt and source hashes, and exact commands.
Paid execution requires an explicit flag and a local cost ceiling. The ceiling
prevents accidental launches but is not a provider-side billing limit.

Run outputs archive structured predictions, raw responses, token usage, model
metadata, errors, and timing. Metrics are regenerated from matching gold and
prediction manifests rather than copied from historical summary files. The
curated thesis package contains aggregate metrics and frozen configuration but
excludes source screenshots and per-item raw interactions.

Automated tests cover schemas, aggregation gates, evaluation metrics, run
preparation, stability analysis, open-model adapters, and package auditing. Test
success supports implementation consistency; it does not substitute for
empirical validation of labels or evidence.

\section{Implementation Boundaries}\label{implementation-boundaries}

The system is a research prototype. Lexical retrieval is reproducible but not a
strong semantic retriever. OCR quality varies with resolution and interface
style. Hosted model behavior may change despite fixed identifiers, and
commercial APIs do not expose their complete serving stack. The local
open-weight pilot was too slow for a complete benchmark on the available
hardware, so the open-weight comparison uses hosted inference and reports this
limitation.

The architecture supports uploaded screenshot flows, but the thesis benchmark
is fixed to the reviewed Mind2Web-derived set. The implementation can display
region-level evidence even where a final human grounding evaluation is not yet
complete. Product functionality and validated scientific contribution are
therefore kept distinct throughout the following chapters.
